\documentclass[tikz,border=10pt]{standalone}
\usepackage{tikz}
\usetikzlibrary{positioning, arrows.meta, 3d, calc, shapes, backgrounds, fit, shadows}

% Define Colors
\definecolor{techblue}{RGB}{50, 100, 200}
\definecolor{lightblue}{RGB}{180, 210, 255}
\definecolor{warmorange}{RGB}{255, 140, 50}
\definecolor{coralred}{RGB}{255, 100, 80}
\definecolor{fusionpurple}{RGB}{140, 50, 200}
\definecolor{emeraldgreen}{RGB}{40, 180, 100}

\begin{document}
\begin{tikzpicture}[
    x={(1cm,0cm)}, y={(0cm,1cm)}, z={(0.4cm,0.3cm)},
    font=\sffamily\footnotesize,
    >=Stealth,
    node distance=1.5cm
]

% Styles
\tikzstyle{imgnode} = [draw=black!30, thick, inner sep=0pt, fill=white, drop shadow]
\tikzstyle{trapezoid_backbone} = [
    top color=white, bottom color=black!10,
    draw=black!50, thick,
    shape=trapezium, trapezium angle=75,
    minimum height=1.2cm, minimum width=2.5cm,
    trapezium stretches=true,
    inner sep=0pt
]
\tikzstyle{cube} = [
    draw=black!60, line width=0.5pt,
    fill opacity=0.8
]
\tikzstyle{vector} = [
    draw=black!60, rounded corners=2pt,
    line width=0.5pt
]
\tikzstyle{op} = [
    circle, draw=black, thick, fill=white, inner sep=0pt, minimum size=0.4cm, font=\large
]

% =================================================================================
% PART 1: INPUTS
% =================================================================================

% --- Top Branch (RGB) ---
\node[imgnode] (input_rgb) at (0, 3, 0) {
    \includegraphics[width=2cm, height=2cm]{rgb_input.jpg}
};
\node[below=0.2cm of input_rgb] {Static RGB Frame (Apex)};

% --- Bottom Branch (Dynamic) ---
\node[imgnode] (input_dyn) at (0, -3, 0) {
    \includegraphics[width=2cm, height=2cm]{dynamic_input.jpg}
};
\node[below=0.2cm of input_dyn] {Dynamic Image (Motion)};


% =================================================================================
% PART 2: BACKBONES (ResNet-18)
% =================================================================================

% --- RGB Backbone ---
% 3D Trapezoid Effect manually drawing to look better than shape=trapezium in 3D
\coordinate (bb_rgb_start) at (2.5, 3, 0);
\coordinate (bb_rgb_end) at (5.5, 3, 0);

% Front Face (Trapezoid)
\draw[fill=techblue!20, draw=techblue!80, thick] 
    ($(bb_rgb_start)+(0, 0.8, 0)$) coordinate (rtl) -- 
    ($(bb_rgb_start)+(0, -0.8, 0)$) coordinate (rbl) -- 
    ($(bb_rgb_end)+(0, -0.4, 0)$) coordinate (rbr) -- 
    ($(bb_rgb_end)+(0, 0.4, 0)$) coordinate (rtr) -- cycle;
% Side/Top (3D effect)
\draw[fill=techblue!10, draw=techblue!80, thick] (rtr) -- (rbr) -- ($(rbr)+(0,0,-1)$) -- ($(rtr)+(0,0,-1)$) -- cycle;
\draw[fill=techblue!10, draw=techblue!80, thick] (rtl) -- (rtr) -- ($(rtr)+(0,0,-1)$) -- ($(rtl)+(0,0,-1)$) -- cycle;

\node[align=center, text=techblue!80!black] at (4, 3, 0) {ResNet-18\\Backbone};

% Output Feature Map (RGB) - Cube
\coordinate (feat_rgb_pos) at (6.5, 3, 0);
\draw[cube, fill=techblue!40] (feat_rgb_pos) -- ++(0, 0.8, 0) -- ++(0.8, 0, 0) -- ++(0, -0.8, 0) -- cycle;
\draw[cube, fill=techblue!30] ($(feat_rgb_pos)+(0,0.8,0)$) -- ++(0,0,-0.8) -- ++(0.8,0,0) -- ++(0,0,0.8) -- cycle;
\draw[cube, fill=techblue!50] ($(feat_rgb_pos)+(0.8,0,0)$) -- ++(0,0.8,0) -- ++(0,0,-0.8) -- ++(0,-0.8,0) -- cycle;
\node[below=0.5cm of feat_rgb_pos, xshift=0.4cm] {$C \times H \times W$};


% --- Dynamic Backbone ---
\coordinate (bb_dyn_start) at (2.5, -3, 0);
\coordinate (bb_dyn_end) at (5.5, -3, 0);

% Front Face
\draw[fill=warmorange!20, draw=warmorange!80, thick] 
    ($(bb_dyn_start)+(0, 0.8, 0)$) coordinate (dtl) -- 
    ($(bb_dyn_start)+(0, -0.8, 0)$) coordinate (dbl) -- 
    ($(bb_dyn_end)+(0, -0.4, 0)$) coordinate (dbr) -- 
    ($(bb_dyn_end)+(0, 0.4, 0)$) coordinate (dtr) -- cycle;
% 3D sides
\draw[fill=warmorange!10, draw=warmorange!80, thick] (dtr) -- (dbr) -- ($(dbr)+(0,0,-1)$) -- ($(dtr)+(0,0,-1)$) -- cycle;
\draw[fill=warmorange!10, draw=warmorange!80, thick] (dtl) -- (dtr) -- ($(dtr)+(0,0,-1)$) -- ($(dtl)+(0,0,-1)$) -- cycle;

\node[align=center, text=warmorange!80!black] at (4, -3, 0) {ResNet-18\\Backbone};

% Output Feature Map (Dynamic) - Cube
\coordinate (feat_dyn_pos) at (6.5, -3, 0);
\draw[cube, fill=warmorange!40] (feat_dyn_pos) -- ++(0, 0.8, 0) -- ++(0.8, 0, 0) -- ++(0, -0.8, 0) -- cycle;
\draw[cube, fill=warmorange!30] ($(feat_dyn_pos)+(0,0.8,0)$) -- ++(0,0,-0.8) -- ++(0.8,0,0) -- ++(0,0,0.8) -- cycle;
\draw[cube, fill=warmorange!50] ($(feat_dyn_pos)+(0.8,0,0)$) -- ++(0,0.8,0) -- ++(0,0,-0.8) -- ++(0,-0.8,0) -- cycle;
\node[below=0.5cm of feat_dyn_pos, xshift=0.4cm] {$C \times H \times W$};


% Arrows from Input to Backbone to Feature
\draw[->, thick, techblue] (input_rgb) -- (2.5, 3, 0);
\draw[->, thick, techblue] (5.5, 3, 0) -- (6.5, 3, 0);
\draw[->, thick, warmorange] (input_dyn) -- (2.5, -3, 0);
\draw[->, thick, warmorange] (5.5, -3, 0) -- (6.5, -3, 0);


% =================================================================================
% PART 3: ASYMMETRIC MODULES
% =================================================================================

% --- TOP: CBAM + Grid Pooling ---

% 1. CBAM Channel
\coordinate (cbam_start) at (8.0, 3, 0);
\draw[->, thick] ($(feat_rgb_pos)+(0.8, 0.4, 0)$) -- (cbam_start);

% Channel Slice
\draw[cube, fill=red!40] (cbam_start) -- ++(0, 0.8, 0) -- ++(0.2, 0, 0) -- ++(0, -0.8, 0) -- cycle;
\draw[cube, fill=red!30] ($(cbam_start)+(0,0.8,0)$) -- ++(0,0,-0.8) -- ++(0.2,0,0) -- ++(0,0,0.8) -- cycle;
\draw[cube, fill=red!50] ($(cbam_start)+(0.2,0,0)$) -- ++(0,0.8,0) -- ++(0,0,-0.8) -- ++(0,-0.8,0) -- cycle;

\node[op, scale=0.7] (mul1) at (8.8, 3.4, 0) {$\otimes$};

% 2. CBAM Spatial
\coordinate (spatial_start) at (9.5, 3, 0);
\draw[->, thick] (mul1) -- (spatial_start);
% Spatial Grid (Flat)
\draw[fill=gray!20, draw=black, step=0.2] (spatial_start) grid ++(0.8, 0.8);
\node[op, scale=0.7] (mul2) at (10.8, 3.4, 0) {$\otimes$};

\draw[->, thick] ($(spatial_start)+(0.8, 0.4)$) -- (mul2);

\node[below=0.6cm of spatial_start, align=center, font=\scriptsize] {CBAM\\(Attention)};

% 3. Grid Pooling
\coordinate (grid_start) at (11.8, 2.6, 0); % Lowered slightly to center
\draw[->, thick] (mul2) -- ($(grid_start)+(-0.2, 0.8, 0)$);

% 7x7 Grid (Symbolic)
\draw[fill=lightblue!20, draw=techblue, step=0.4] (grid_start) grid ++(1.6, 1.6);
% Bold lines for 2x2 split
\draw[techblue!80!black, very thick] ($(grid_start)+(0.8, 0)$) -- ++(0, 1.6);
\draw[techblue!80!black, very thick] ($(grid_start)+(0, 0.8)$) -- ++(1.6, 0);

% Arrows converging
\coordinate (vec_rgb_start) at (14.5, 3.4, 0);
\draw[->, techblue] ($(grid_start)+(0.4, 1.2)$) to[out=0,in=180] (vec_rgb_start);
\draw[->, techblue] ($(grid_start)+(1.2, 1.2)$) to[out=0,in=180] (vec_rgb_start);
\draw[->, techblue] ($(grid_start)+(0.4, 0.4)$) to[out=0,in=180] (vec_rgb_start);
\draw[->, techblue] ($(grid_start)+(1.2, 0.4)$) to[out=0,in=180] (vec_rgb_start);

\node[below=0.2cm of grid_start, font=\scriptsize] {Grid Pooling (2x2)};

% 4. 1D Vector RGB
\draw[vector, fill=techblue!60] (vec_rgb_start) rectangle ++(2.0, 0.4);
\node[above=0.1cm of vec_rgb_start, xshift=1cm, font=\scriptsize, text=techblue!80!black] {Spatial Vector};


% --- BOTTOM: SPP ---

\coordinate (spp_start) at (8.5, -3, 0);
\draw[->, thick] ($(feat_dyn_pos)+(0.8, 0.4, 0)$) -- (spp_start);

% Level 1: 1x1
\draw[fill=warmorange!20, draw=warmorange] (spp_start) rectangle ++(1.0, 1.0);
\node at ($(spp_start)+(0.5, 0.5)$) {1x1};

% Level 2: 2x2
\coordinate (spp_l2) at ($(spp_start)+(0, -1.2)$);
\draw[fill=warmorange!20, draw=warmorange, step=0.5] (spp_l2) grid ++(1.0, 1.0);
\node at ($(spp_l2)+(0.5, 0.5)$) {2x2};

% Level 3: 4x4
\coordinate (spp_l3) at ($(spp_l2)+(0, -1.2)$);
\draw[fill=warmorange!20, draw=warmorange, step=0.25] (spp_l3) grid ++(1.0, 1.0);
\node at ($(spp_l3)+(0.5, 0.5)$) {4x4};

\node[below=0.2cm of spp_l3, font=\scriptsize] {Spatial Pyramid Pooling (SPP)};

% Merging Arrows
\coordinate (vec_dyn_start) at (12.5, -3, 0); % Shifted left to make vector longer
\draw[->, warmorange] ($(spp_start)+(1.0, 0.5)$) -- ($(vec_dyn_start)+(0, 0.2)$);
\draw[->, warmorange] ($(spp_l2)+(1.0, 0.5)$) -- ($(vec_dyn_start)+(0, 0.2)$);
\draw[->, warmorange] ($(spp_l3)+(1.0, 0.5)$) -- ($(vec_dyn_start)+(0, 0.2)$);

% Super Long 1D Vector Dynamic
\draw[vector, fill=warmorange!60] (vec_dyn_start) rectangle ++(4.0, 0.4);
\node[below=0.1cm of vec_dyn_start, xshift=2cm, font=\scriptsize, text=warmorange!80!black] {Temporal Vector};


% =================================================================================
% PART 4: FUSION & SE-BLOCK
% =================================================================================

% Concat Point
\coordinate (concat_x) at (17.5, 0, 0);

% Routing lines to concat
\draw[->, thick, techblue!80] ($(vec_rgb_start)+(2.0, 0.2)$) -| (concat_x);
\draw[->, thick, warmorange!80] ($(vec_dyn_start)+(4.0, 0.2)$) -| (concat_x);

% Concatenated Vector
\coordinate (vec_fused_start) at (18.0, -0.5, 0);
% Top half blue
\draw[vector, fill=techblue!60] (vec_fused_start) rectangle ++(2.5, 0.5);
% Bottom half orange
\draw[vector, fill=warmorange!60] ($(vec_fused_start)+(0, -0.5)$) rectangle ++(2.5, 0.5);

\node[above=0.6cm of vec_fused_start, xshift=1.25cm, font=\bfseries] {Concat};

% --- SE Block (U-Shape) ---
\coordinate (se_start) at (19.25, -1.0, 0); % Bottom of fused vector
\coordinate (se_bottom) at (19.25, -4.5, 0);

% Down arrow
\draw[->, thick, fusionpurple] (se_start) -- (se_bottom);

% SE Components along the bottom
% GAP
\node[draw, fill=white, minimum size=0.6cm, font=\scriptsize] (gap) at ($(se_bottom)+(1.0, 0)$) {GAP};
\draw[->, thick, fusionpurple] (se_bottom) -- (gap);

% FC1
\node[draw, fill=fusionpurple!20, minimum height=0.4cm, minimum width=0.2cm] (fc1) at ($(gap)+(1.2, 0)$) {};
\node[below=0.1cm of fc1, font=\tiny] {FC};
\draw[->, thick, fusionpurple] (gap) -- (fc1);

% ReLU
\node[circle, draw=fusionpurple, inner sep=1pt, font=\scriptsize] (relu) at ($(fc1)+(0.8, 0)$) {$\delta$};
\draw[->, thick, fusionpurple] (fc1) -- (relu);

% FC2
\node[draw, fill=fusionpurple!20, minimum height=0.8cm, minimum width=0.2cm] (fc2) at ($(relu)+(0.8, 0)$) {};
\node[below=0.1cm of fc2, font=\tiny] {FC};
\draw[->, thick, fusionpurple] (relu) -- (fc2);

% Sigmoid
\node[circle, draw=fusionpurple, inner sep=1pt, font=\scriptsize] (sig) at ($(fc2)+(0.8, 0)$) {$\sigma$};
\draw[->, thick, fusionpurple] (fc2) -- (sig);

% Weight Vector (Gradient)
\coordinate (w_vec) at ($(sig)+(1.0, 0)$);
\shade[left color=white, right color=fusionpurple] (w_vec) rectangle ++(0.4, 1.0);
\draw[black!50] (w_vec) rectangle ++(0.4, 1.0);
\node[below=0.1cm of w_vec, xshift=0.2cm] {$w$};

\draw[->, thick, fusionpurple] (sig) -- (w_vec);

% Up arrow to main path
\coordinate (mul_se) at (22.5, -0.5, 0);
\draw[->, thick, fusionpurple] ($(w_vec)+(0.2, 1.0)$) |- (mul_se);

% Multiplication Node on Main Path
\draw[dashed, thick] ($(vec_fused_start)+(2.5, 0)$) -- (mul_se);
\node[op, draw=fusionpurple, thick] (op_se) at (mul_se) {$\otimes$};

% Dotted Box for SE
\node[draw=fusionpurple, dashed, fit=(gap) (sig) (w_vec) (se_bottom), inner sep=0.3cm, label={[fusionpurple]below:SE-Attention}] {};


% =================================================================================
% PART 5: CLASSIFICATION
% =================================================================================

% Fused Feature Vector
\coordinate (final_vec) at (24.0, -0.5, 0);
\draw[->, thick] (op_se) -- (final_vec);
\draw[vector, left color=techblue!40, right color=warmorange!40] (final_vec) rectangle ++(1.5, 1.0);
\node[above=0.1cm of final_vec, xshift=0.75cm, font=\scriptsize] {Fused Feature};

% FC & Dropout
\coordinate (cls_block) at (26.5, 0, 0);
\draw[->, thick] ($(final_vec)+(1.5, 0.5)$) -- (cls_block);
\node[draw, fill=emeraldgreen!20, minimum height=1.5cm, minimum width=1.0cm, align=center, font=\scriptsize] (fc_drop) at (27.0, 0, 0) {FC \&\\Dropout};

% Output Classes
\coordinate (out_pos) at (29.0, 0, 0);
\node[circle, fill=emeraldgreen!40, draw=emeraldgreen!80, minimum size=0.8cm] (pos) at (29.0, 1.5, 0) {Pos};
\node[circle, fill=emeraldgreen!40, draw=emeraldgreen!80, minimum size=0.8cm] (neg) at (29.0, 0, 0) {Neg};
\node[circle, fill=emeraldgreen!40, draw=emeraldgreen!80, minimum size=0.8cm] (sur) at (29.0, -1.5, 0) {Sur};

\draw[->, thick, emeraldgreen] (fc_drop.east) -- (pos.west);
\draw[->, thick, emeraldgreen] (fc_drop.east) -- (neg.west);
\draw[->, thick, emeraldgreen] (fc_drop.east) -- (sur.west);

\end{tikzpicture}
\end{document}
